\section{Conclusão}

Nesse trabalho, foi realizado um experimento utilizando o aparato de Millikan, 
onde foi analizado o comportamento de gotículas eletrizadas sobre o efeito de um 
campo elétrico alternante. 
Medindo a velocidade com que as gotas sobem e descem, 
foi possível calcular o valor do raio $r$ e da carga $q$ de cada gota. 

Com os dados obtidos, foi possível observar que $q$ não é função de $r$. 
Além disso, foi oservada uma tendência dos valores de $q$ a se agrupar 
em torno de múltiplos inteiros de $\qty{1,602e-19}{\coulomb}$, 
como demonstrado por Millikan em 1909. 
No entanto, apenas os dados obtidos não são suficientes para considerar 
demonstrada a quantização da carga elétrica, 
especialmente pela falta de pontos com $q$ alto.
Faz-se então necessário a coleta de mais dados, e com mais mecanismos para 
amenizar os desvios experimentais. 